\documentclass[11pt]{article}
\usepackage[margin=1in]{geometry}
\usepackage[T1]{fontenc}
\usepackage{amsmath,amssymb,amsthm,booktabs,hyperref,microtype}
\hypersetup{colorlinks=true,linkcolor=blue,urlcolor=blue,citecolor=blue,
 pdftitle={Sharp asymptotic separation of opposing egalitarian core selections},
 pdfauthor={Maxwell Flanagan}}
\newtheorem{theorem}{Theorem}
\newtheorem{lemma}[theorem]{Lemma}
\newtheorem{proposition}[theorem]{Proposition}
\newcommand{\Lmin}{\operatorname{Lmin}}
\newcommand{\Lmax}{\operatorname{Lmax}}
\newcommand{\Core}{\operatorname{Core}}
\newcommand{\one}{\mathbf1}
\title{Sharp asymptotic separation of opposing\\egalitarian core selections}
\author{Maxwell Flanagan}
\date{September 5, 2026 --- Draft for specialist review}
\begin{document}
\maketitle

\begin{abstract}
For singleton-normalized exact transferable-utility games on $n$ players,
let $s_n$ denote the supremum of the normalized $\ell_1$ distance between
the leximin core allocation and the core allocation that lexicographically
minimizes decreasingly sorted payoffs. We prove
\[
 \lim_{n\to\infty}n(2-s_n)=2,
 \qquad 2-s_n=\frac2n+O\!\left(\frac{\log n}{n^2}\right).
\]
The universal bound holds for opposing maximin and minimax points in every
compact convex nonnegative constant-total set. For the matching asymptotic
construction, we realize a line segment as the entire core of the minimum
of two additive games. Complementary coalition equalities encode a
Fibonacci recurrence. Both full selectors are uniquely determined at their
first optimization stages. The construction exists for every sufficiently
large $n$ and admits integer generators with polynomially bounded total.
The lower-envelope framework, segment-core examples and Fibonacci matrix
ingredients have earlier precedents; the proposed contribution is their
quantitative opposing-selector comparison.
The argument is self-contained; finite exact computations are supplied as
reproducibility checks.
\end{abstract}

\section{The question and the main result}

The core imposes coalition stability but often leaves several efficient
allocations available. Two egalitarian refinements compare those allocations
in opposite orders: leximin first improves the poorest payoff, whereas the
opposing selection first lowers the richest payoff. How far apart can their
recommendations be, even when every coalition value is attained somewhere
in the core?

Let $N$ be a finite player set, $n=|N|$, and let a TU game be a function
$v:2^N\to\mathbb R$ with $v(\varnothing)=0$. Write
$y(S)=\sum_{i\in S}y_i$ and
\[
 \Core(v)=\{y\in\mathbb R^N:y(N)=v(N),\ y(S)\ge v(S)
                                      \text{ for all }S\subseteq N\}.
\]
The game is \emph{exact} if its core is nonempty and, for every coalition
$S$, some core allocation $y$ satisfies $y(S)=v(S)$. Throughout we impose
\emph{singleton normalization}, $v(\{i\})=0$ for all $i$, and a positive
grand total $T=v(N)$. Its core is then a compact convex subset of
$\{y\ge0:y(N)=T\}$.

For $y\in\mathbb R^N$, denote its increasingly and decreasingly sorted
vectors by $y^\uparrow$ and $y^\downarrow$. We use the following explicit
conventions:
\[
 x=\Lmin(v):\quad x^\uparrow
       \text{ is lexicographically maximal over }\Core(v),
\]
\[
 z=\Lmax(v):\quad z^\downarrow
       \text{ is lexicographically minimal over }\Core(v).
\]
Thus $\Lmax$ means the lexicographic refinement of \emph{minimizing} the
maximum payoff. Both rules sort all individual payoffs. Neither rule here
sorts coalition excesses.

These full selections exist and are unique on a nonempty compact convex
set. Existence follows by successive optimization of continuous order
statistics on compact sets. For uniqueness of the first rule, two distinct
optimal points would have the same sorted vector. Consider the smallest
payoff level whose assigned player sets differ. Their midpoint keeps all
lower levels unchanged and has fewer coordinates at this level, strictly
improving the sorted vector. Convexity gives a contradiction. Applying the
same argument to the negative of the feasible set proves uniqueness of the
second rule.

Put
\begin{equation}\label{eq:distance}
 O(x,z)=\sum_{i\in N}\min(x_i,z_i),\qquad
 \rho(v)=\frac{\|x-z\|_1}{T}=2-\frac{2O(x,z)}T,
 \qquad s_n=\sup_v\rho(v),
\end{equation}
where the supremum ranges over all singleton-normalized exact games with
$n$ players and positive total. Scaling the game by a positive constant
leaves $\rho$ unchanged. No attainment of this supremum is assumed.

\begin{theorem}\label{thm:main}
For every $n\ge2$,
\[
 s_n\le2-\frac2{n-1}.
\]
For every integer $n\ge128$, let $L=\lceil\log_2 n\rceil$. Then
\begin{equation}\label{eq:finite-squeeze}
 \frac2{n-1}\le 2-s_n
       \le\frac2{n-12L}+\frac2{n^2+1}.
\end{equation}
Consequently
\begin{equation}\label{eq:asymptotic}
 \lim_{n\to\infty}n(2-s_n)=2,
 \qquad 2-s_n=\frac2n+O\!\left(\frac{\log n}{n^2}\right).
\end{equation}
The lower construction for $s_n$ uses two nonnegative efficient generators
and has a one-dimensional entire core. It can be scaled to an integer-valued
game; an alternative parameter choice gives integer generator coordinates
of $O(\log n)$ bits and total $O(n^3)$.
\end{theorem}

The normalized vectors $x/T$ and $z/T$ are probability vectors, so
$\rho/2$ is their total variation distance. The result concerns disagreement
about distributing a fixed total: it does not assert a loss of aggregate
worth. Fixed-$n$ sharpness, an exact second-order term, and uniqueness of
asymptotic extremizers are not asserted.

\paragraph{Relation to prior work.}
The leximin stable allocation in the core is established in Arin and
I\~narra~\cite{arin-inarra}. Vazirani~\cite{vazirani}, notably Sections 1.2
and 5, studies the two opposing fairness criteria for assignment-game
cores and gives efficient algorithms; his Lemma 3 also proves uniqueness
of leximin. Gangam et al.~\cite{gangam} study corresponding algorithms for
max-flow, MST, and $b$-matching games, with their constructive algorithms
operating on Owen sets. These domains and algorithmic objectives differ
from the unrestricted exact-game distance supremum considered here.
General max-min/min-max fairness and convex-domain uniqueness also predate
these applications \cite{radunovic,abernethy}. Price-of-fairness bounds
\cite{bertsimas} concern aggregate utility loss, whereas our metric compares
recipients of a fixed total. Individual-ideal normalization in that framework
can also change the selector; its bounds do not directly imply ours.

Exactness and lower-envelope representations are classical. For example,
Amarante et al.~\cite{amarante}, Section 2.2, record both the lower-envelope
interpretation and the fact that its core may strictly contain the
generating convex hull. Studen\'y and Kratochv\'il~\cite{studeny} give a
facet description of the finite exact-game cone. Our proof uses neither a
facet enumeration nor an unproved replacement of the core by its generators.
The minimum of two equal-total additive proposals is already the
\emph{bifocal distribution game} \cite{bifocal}, and classical examples
already realize entire segment cores \cite{shellshear}. Neither device is
claimed as new.

The Fibonacci ingredient also has a concrete antecedent. The weighted-majority
family attributed to Isbell by Pressacco and Ziani \cite{pressacco15}, for
$p\ge4$, has
$p$ positive weights
$(F_1,\ldots,F_{p-2},F_{p-2},F_{p-1})$ and quota $F_p$.
The minimal-winning-coalition matrix is invertible by
\cite[Theorem~4.2]{pressacco14}. Dividing the weights by $F_p$ therefore
supplies a positive vector $d=A^{-1}\one$ with
$\min_i d_i=1/F_p$ and $\one^Td=2-1/F_p$.
Our inference is that this older ingredient can be substituted into
Lemma~\ref{lem:segment} to obtain the same asymptotic order with
$p=O(\log n)$ donor coordinates. Accordingly the matrix construction is a
specialization of established ingredients, not a new general Fibonacci
method. We retain the self-contained construction below for reproducibility.

The proposed contribution is the universal overlap bound and the sharp
first-order distance theorem. A bounded source comparison has not identified
an equivalent quantitative result; this does not certify novelty.
Specialist priority and significance assessment remains pending. The
companion source ledger records the versions, deductions and access gaps.

\section{A universal overlap inequality}

\begin{lemma}\label{lem:overlap}
Let $n\ge2$, $T>0$, and let $K$ be a nonempty compact convex subset of
$\{y\in\mathbb R^n:y\ge0,\ y(N)=T\}$. If
\[
 x\in\arg\max_{y\in K}\min_i y_i,
 \qquad z\in\arg\min_{y\in K}\max_i y_i,
\]
then $\sum_i\min(x_i,z_i)\ge T/(n-1)$.
\end{lemma}

\begin{proof}
If $x=z$, the overlap is $T$. Otherwise let $a=\min_i x_i$ and
$M=\max_i z_i$. Since $a\le T/n\le M$, equality $a=M$ would force
both points to be uniform; hence $a<M$.

Some index $i$ with $x_i=a$ satisfies $z_i\le a$. If every minimum
coordinate instead had $z_i>a$, a sufficiently small positive move from
$x$ toward $z$ would raise all of them, while all other coordinates would
remain above $a$. This contradicts maximin optimality. Similarly some
index $j$ with $z_j=M$ has $x_j\ge M$: otherwise a small move from $z$
toward $x$ lowers all maximum coordinates and keeps the others below $M$.
The two indices are distinct because $x_i=a<M\le x_j$. Therefore
\[
 O(x,z)\ge z_i+M,
\]
and nonnegativity gives
\[
 T\le z_i+(n-1)M
   \le(n-1)(z_i+M)
   \le(n-1)O(x,z).
\]
\end{proof}

By \eqref{eq:distance}, the lemma proves the upper bound on $s_n$ in
Theorem~\ref{thm:main}. Exactness and all later lexicographic stages are
unnecessary for this bound.

\section{Realizing the entire core as a segment}

\begin{lemma}\label{lem:segment}
Let $p\ge1$ and let $A\in\{0,1\}^{p\times p}$ be invertible with
$d=A^{-1}\one>0$. Choose an integer $r\ge1$ such that
$H=r-\sum_i d_i\ge1$. On one distinguished donor, $p$ other donors, and
$r$ receivers, define
\[
 U=(H,d_1,\ldots,d_p;0^r),\qquad
 V=(0^{p+1};1^r),\qquad
 v(S)=\min\{U(S),V(S)\}.
\]
Then $v$ is singleton-normalized and exact, has total $r$, and
$\Core(v)=[U,V]$. If $\alpha=\min_i d_i$, its full selectors equal
$tU+(1-t)V$ at
\begin{equation}\label{eq:selector-t}
 t_{\min}=\frac1{1+\alpha},\qquad
 t_{\max}=\frac1{1+H}.
\end{equation}
In particular,
\begin{equation}\label{eq:segment-deficit}
 2-\rho(v)=2\left(\frac\alpha{1+\alpha}+\frac1{1+H}\right).
\end{equation}
\end{lemma}

\begin{proof}
Both generators are nonnegative, have total $r$, and satisfy all coalition
lower bounds. Every coalition value is attained by one of them, establishing
exactness. Their supports are disjoint, establishing singleton normalization.
Since every column of $A$ is nonzero, positivity and $Ad=\one$ imply
$d_i\le1$ for every $i$.

Let $D$ be the subset of nondistinguished donors specified by any row of
$A$, and let $e$ be any receiver. Both generators give $D\cup\{e\}$
value one, and give its complement value $r-1$. Thus efficiency and these
two coalition inequalities force, at every core point,
\[
 y(D)+y_e=1.
\]
Fixing a row and varying $e$ shows that all receivers have a common
coordinate, say $1-t$. Varying the row gives $Ay_d=t\one$, whence
$y_d=td$. Efficiency gives $y_h=tH$ for the distinguished donor.
Nonnegativity now implies $0\le t\le1$. Conversely every point of
$[U,V]$ is in the core by convexity. This proves the whole-core equality.

On the segment the smallest and largest coordinates are
\[
 \min\{t\alpha,1-t\},\qquad \max\{tH,1-t\},
\]
because $0<\alpha\le d_i\le1\le H$. Each has a unique relevant scalar
optimum, giving \eqref{eq:selector-t}. These unique first-stage optima
determine the full lexicographic selectors. Since
$\|U-V\|_1=2r$ and $t_{\min}\ge t_{\max}$,
\[
 \rho(v)=2(t_{\min}-t_{\max}),
\]
which rearranges to \eqref{eq:segment-deficit}.
\end{proof}

The essential restriction is the system of unit-weight donor subsets.
Merely defining a game as a lower envelope proves exactness but does not
generally prove that its core is the generating segment.

\section{The Fibonacci construction and all player counts}

Let $F_1=F_2=1$ and $F_{j+2}=F_{j+1}+F_j$. For integers $m\ge2$ and
$r\ge m+1$, set
\[
 F=F_{m+1},\quad H=r-m,\quad
 a_i=F_{m-i+1}/F,\quad b_i=1-a_i\quad(1\le i\le m).
\]
Use donor labels $h,A_1,\ldots,A_m,B_1,\ldots,B_m$ and $r$ receivers,
and define
\begin{equation}\label{eq:generators}
 U=(H,a_1,\ldots,a_m,b_1,\ldots,b_m;0^r),\qquad
 V=(0^{2m+1};1^r).
\end{equation}
Thus $n=r+2m+1$ and $T=r$.

\begin{proposition}\label{prop:fibonacci}
The game $v(S)=\min\{U(S),V(S)\}$ from \eqref{eq:generators} is
singleton-normalized and exact, its entire core is $[U,V]$, and
\begin{equation}\label{eq:fibonacci-deficit}
 \rho(v)=2\left(\frac{F}{F+1}-\frac1{r-m+1}\right),
 \qquad 2-\rho(v)=\frac2{F_{m+1}+1}+\frac2{r-m+1}.
\end{equation}
\end{proposition}

\begin{proof}
The following $2m$ subsets of the nondistinguished donors have $U$-weight
one:
\begin{align*}
 &\{A_i,B_i\} &&(1\le i\le m),\\
 &\{A_{i+1},A_{i+2},B_i\} &&(1\le i\le m-2),\\
 &\{A_{m-1},B_m\},\qquad \{A_1,A_2\}.&&
\end{align*}
The pair rows are immediate; the other equalities are the recurrence and
the identities $a_{m-1}=a_m$ and $a_1+a_2=1$.

Their incidence matrix is invertible. In its homogeneous system the pair
rows eliminate $B_i=-A_i$. The other rows then give
\[
 A_i=A_{i+1}+A_{i+2},\qquad A_{m-1}=A_m,
 \qquad A_1+A_2=0.
\]
Writing $A_m=u$ gives $A_i=F_{m-i+1}u$ and $F_{m+1}u=0$, so every
coordinate vanishes. This also covers $m=2$, when recurrence rows are
absent. The positive inhomogeneous solution $(a,b)$ has total $m$ and
minimum $1/F$: the smallest $a_i$ is $1/F$, while
$\min_i b_i=F_{m-1}/F\ge1/F$. Apply Lemma~\ref{lem:segment} with
$p=2m$, $\alpha=1/F$, and $H=r-m$.
\end{proof}

\begin{proof}[Proof of Theorem~\ref{thm:main}]
Lemma~\ref{lem:overlap} gives the left inequality in
\eqref{eq:finite-squeeze}. For its right inequality choose, for every
$n\ge128$,
\[
 L=\lceil\log_2 n\rceil,\qquad m=4L,\qquad r=n-8L-1.
\]
We need $r\ge m+1$, equivalently $n\ge12L+2$. At $n=128$ this is a
direct check. For $L\ge8$, the least possible $n$ is $2^{L-1}+1$;
$2^{L-1}+1\ge12L+2$ holds at $L=8$ and its left side grows faster
thereafter. Also $F_{2j+1}\ge2^j$ for $j\ge0$, by induction from
$F_{2j+3}\ge2F_{2j+1}$. Hence
\[
 F_{m+1}\ge2^{2L}\ge n^2,\qquad r-m+1=n-12L>0.
\]
Proposition~\ref{prop:fibonacci} proves the required upper bound on the
deficit for this exact game. Taking the supremum of separation gives
\eqref{eq:finite-squeeze}. Finally,
\[
 \frac2{n-12L}-\frac2n=\frac{24L}{n(n-12L)},
\]
and $L=O(\log n)$, proving \eqref{eq:asymptotic} by squeezing.

Multiplication by $F$ makes every generator coordinate integral. For the
polynomial-total claim choose instead the least $m\ge2$ with
$F_{m+1}\ge n^2$ and again put $r=n-2m-1$. For sufficiently large $n$,
$m=O(\log n)$, $r\ge m+1$, and
$n^2\le F_{m+1}<2n^2$. The scaled total is $rF_{m+1}=O(n^3)$, and
every nonnegative coordinate is bounded by that total. Thus each coordinate
uses $O(\log n)$ bits. The two vectors are a compact representation of
the game; listing all coalition values would still require exponentially
many entries.
\end{proof}

\section{Verification and scope of the evidence}

The proof above does not depend on a computational census.
The companion code reconstructs every proper coalition inequality directly
from the two rational generators. Run \texttt{run\_checks.py} in the
\texttt{verification} directory to check the cases
\[
 (m,r)\in\{(2,3),(2,4),(3,4),(3,5),(4,5)\}.
\]
These cases have $8,9,11,12,14$ players. Exact rational Z3 queries exclude
every off-segment core allocation, compute both first-stage optimum values,
and exclude every other point on either optimal face. Generator feasibility
and convexity establish the opposite core containment algebraically.
Arithmetic checks also cover $2\le m\le64$ and displayed large-$n$
parameter choices.

A deliberate perturbation of one donor pair at $(m,r)=(2,3)$ preserves
the common generator total and singleton normalization but breaks the
unit-sum realization. The verifier finds and checks an explicit core point
outside the proposed segment. This guards against treating every lower
envelope as the convex hull of its generators. The code uses explicit
exceptions, rejects UNKNOWN and timeouts, and has been replayed with and
without Python optimization. Its JSON records are finite corroboration;
they are not a formal proof of the all-parameter theorem. Specialist
correctness and priority review remain pending.

\begin{thebibliography}{99}
\bibitem{radunovic}
B. Radunovi\'c and J.-Y. Le Boudec.
\emph{A Unified Framework for Max-Min and Min-Max Fairness With Applications}.
IEEE/ACM Transactions on Networking 15 (2007), 1073--1083.
\url{https://doi.org/10.1109/TNET.2007.896231}.
\bibitem{abernethy}
J. Abernethy, R. E. Schapire and U. Syed.
\emph{Lexicographic Optimization: Algorithms and Stability}.
arXiv:2405.01387 (2024). \url{https://arxiv.org/abs/2405.01387}.
\bibitem{bertsimas}
D. Bertsimas, V. F. Farias and N. Trichakis.
\emph{The Price of Fairness}. Operations Research 59 (2011), 17--31.
\url{https://www.mit.edu/~nitric/pof_or.pdf}.
\bibitem{bifocal}
P. Gadea-Blanco, J.-M. Gim\'enez-G\'omez and M. C. Marco-Gil.
\emph{Compromising in bifocal distribution games: the average value}.
Theory and Decision 81 (2016), 449--465.
\url{https://doi.org/10.1007/s11238-016-9540-5}.
\bibitem{shellshear}
E. Shellshear and P. Sudh\"olter.
\emph{On core stability, vital coalitions, and extendability}.
Games and Economic Behavior 67 (2009), 633--644.
Full text inspected in Discussion Paper 3/2008, Example 4.3.
\bibitem{pressacco14}
F. Pressacco, G. Plazzotta and L. Ziani.
\emph{Twin relationships in Parsimonious Games: some results}.
arXiv:1404.2305 (2014). \url{https://arxiv.org/abs/1404.2305}.
\bibitem{pressacco15}
F. Pressacco and L. Ziani.
\emph{A Fibonacci Approach to Weighted Majority Games}.
Journal of Game Theory 4 (2015), 36--44.
\url{https://doi.org/10.5923/j.jgt.20150402.03}.
\bibitem{arin-inarra}
J. Arin and E. I\~narra.
\emph{Egalitarian solutions in the core}.
International Journal of Game Theory 30 (2001), 187--193.
\url{https://doi.org/10.1007/s001820100073}.

\bibitem{vazirani}
V. V. Vazirani.
\emph{Fair Core Imputations for the Assignment Game: New Solution Concepts
and Efficient Algorithms}.
arXiv:2402.11437, version 6 (2025).
\url{https://arxiv.org/html/2402.11437v6}.

\bibitem{gangam}
R. R. Gangam, N. Garg, P. Shahkar, and V. V. Vazirani.
\emph{Equitable Core Imputations for Max-Flow, MST and $b$-Matching Games}.
arXiv:2403.06037, version 3 (2025).
\url{https://arxiv.org/pdf/2403.06037v3}.

\bibitem{amarante}
M. Amarante, F. Maccheroni, M. Marinacci, and L. Montrucchio.
\emph{Cores of Non-Atomic Market Games}.
The Carlo Alberto Notebooks 13, May 2006, Section 2.2.
\url{https://www.carloalberto.org/wp-content/uploads/2018/11/no.13.pdf}.

\bibitem{studeny}
M. Studen\'y and V. Kratochv\'il.
\emph{Facets of the cone of exact games}.
Mathematical Methods of Operations Research 95 (2022), 35--80.
\url{https://doi.org/10.1007/s00186-022-00770-4}.
\end{thebibliography}
\end{document}
